% !TEX root = ../main.tex

\chapter{实验脚本与日志归档说明}

本附录结合现有实验脚本与历史运行记录，对论文中各验证阶段所对应的脚本入口、日志归档方式和参数组织口径作统一说明。目的不是列出全部内部文件，而是将正文中需要复核的单卡、多 GPU、真实集有效性和论文出图阶段，整理为可复用的一套归档规范。

\section{实验阶段与材料对应关系}

\small
\begin{tabularx}{\textwidth}{P{0.22\textwidth}P{0.29\textwidth}Y}
论文验证阶段 & 代表性脚本或日志前缀 & 主要用途 \\
真实集路径检查与冒烟验证 & \path|freeze_last_tune_<ts>/01_*| & 对真实集输入路径、基础运行链路和候选召回---评分接口进行快速联调，对应正文中的真实集有效性验证起点。 \\
真实集校准与配置筛选 & \path|freeze_last_tune_<ts>/02_*| & 保存校准轮次的标准输出、排序摘要和配置对比结果，用于冻结论文定稿前的候选规模与评分配置。 \\
单卡链路对比与阶段耗时分析 & \path|bench_flash_<scale>_<scene>_gpu<id>.log|、\path|compare_<topic>_<ts>.md| & 对应第四章单卡路径优化实验，重点记录原始链路、召回加速、GPU 主导链路和完整优化方案之间的耗时与吞吐差异。 \\
多 GPU 扩展、归并与冲突对比 & \path|bench_flash_<scale>_replica_*|、\path|*_shard_*|、\path|gpu_<topic>_gpu<id>.csv| & 对应第五章索引复制主方案、设备侧归并兜底路径、主机侧聚合对比基线以及共享路径冲突实验。 \\
论文图表整理与结果汇总 & \path|latest_<topic>_paper/|、\path|chart_ready_*.csv|、\path|session_meta.txt| & 将多轮实验结果整理为论文图表输入数据，保留轮次汇总、出图口径与会话元信息。 \\
\end{tabularx}
\normalsize

\section{统一文件命名风格}

\begin{itemize}
  \item 会话目录统一采用“任务名 + 时间戳”形式，例如 \path|freeze_last_tune_<YYYYmmdd_HHMMSS>|，用于标识一轮冻结前复核或专题实验。
  \item 单次运行日志统一采用 \path|bench_<route>_<scale>_<scene>_gpu<id>.log|，其中 \texttt{route} 表示单卡或多卡执行路径，\texttt{scale} 表示数据规模，\texttt{scene} 表示验证场景或配置标签。
  \item 对比结果统一采用 \path|compare_<topic>_<timestamp>.md| 与 \path|compare_<topic>_<timestamp>.csv| 成对保存，其中 Markdown 文件用于保留人工可读结论，CSV 文件用于后续绘图与表格汇总。
  \item 过程检查类文件统一采用 \path|0N_<step>.prompt.txt|、\path|0N_<step>.stdout.log| 和 \path|checklist.log| 的组合形式，用于保留命令入口、标准输出和阶段性检查记录。
  \item 出图目录统一保留 \path|all_rounds.csv|、\path|chart_ready_*.csv| 和 \path|session_meta.txt| 三类文件，分别对应原始轮次结果、绘图输入和会话说明。
\end{itemize}

\section{统一参数设置风格}

\begin{itemize}
  \item 数据与划分参数统一使用 \texttt{real-root}、\texttt{real-datasets}、\texttt{synthetic-root}、\texttt{synthetic-scales}、\texttt{input-npz}、\texttt{out-dir}、\texttt{seed} 等名称，明确区分真实集入口、合成集规模和数据切分位置。
  \item 单卡热路径参数统一使用 \texttt{device}、\texttt{batch-size}、\texttt{score-chunk-size}、\texttt{warmup}、\texttt{repeats}、\texttt{max-run-seconds}，分别控制设备绑定、批大小、评分块大小、预热次数和单轮运行上限。
  \item 候选召回与索引参数统一使用 \texttt{flash-search-k}、\texttt{flash-rerank-exact}、\texttt{flash-indices-only}、\texttt{faiss-index}、\texttt{faiss-nlist}、\texttt{faiss-nprobe}、\texttt{faiss-pq-m}、\texttt{faiss-pq-bits} 等名称，避免在不同脚本中混用语义相近的别名。
  \item 多 GPU 相关参数统一使用 \texttt{gpu-list}、\texttt{flash-multi-gpu-route}、\texttt{flash-service-topology}、\texttt{flash-service-protocol}、\texttt{flash-multi-gpu-devices}、\texttt{flash-multi-gpu-gather-device}、\texttt{flash-shards}、\texttt{flash-service-batch-size}，明确区分索引复制主路径、分片兜底路径和归并位置。
  \item 监控与分析参数统一使用 \texttt{gpu-log-path}、\texttt{gpu-log-interval-ms}、\texttt{chain-profile}、\texttt{trace-output-dir} 等名称；Shell 包装脚本中对应使用全大写环境变量，如 \texttt{BATCH\_SIZE}、\texttt{FLASH\_SEARCH\_K}，冻结配置则统一使用 \texttt{FREEZE\_*} 前缀。
\end{itemize}


\chapter{图表补充说明}

本附录对正文图表中反复出现的变量命名、单位和比较口径作统一补充，便于阅读时快速核对不同章节之间的含义是否一致。

\section{图表口径说明}

\begin{itemize}
  \item 图中的 QPS、平均时延和 p99 时延均指对应实验场景下的端到端统计结果；若未特别说明，时延单位统一为 ms。
  \item 单卡结果图主要用于比较数据路径优化前后的链路差异，多 GPU 结果图主要用于比较索引复制主方案、设备侧归并兜底路径以及共享路径冲突条件下的扩展差异。
  \item 真实集结果用于有效性验证，1M 及以上合成数据结果用于系统性能与扩展分析，二者不直接比较绝对吞吐数值。
  \item 正文结果图中的方案名称均对应具体验证场景或系统路径，不再使用含义不明确的编号标签。
\end{itemize}

\section{特征构造字段与派生特征补充说明}

\begin{table}[htbp]
  \caption{特征构造明细}
  \label{tab:feature-engineering}
  \centering
  \small
  \renewcommand{\arraystretch}{1.2}
  \begin{adjustbox}{max width=\linewidth}
  \begin{tabular}{P{2.8cm}P{4.1cm}P{2.4cm}P{3.2cm}P{4.3cm}}
    \toprule
    原始字段/信息源 & 派生特征 & 特征类型 & 预处理方式 & 在系统中的作用\\
    \midrule
    交易金额 & 标准化金额、对数金额、时间窗金额均值/波动统计 & 连续型 & 缺失填补、异常值截断、标准化/归一化 & 刻画交易规模及金额异常波动\\
    交易时间戳 & 小时段、星期信息、夜间交易标记、平均交易间隔 & 时间型/统计型 & 时间规范化、周期化表达、滑动时间窗统计 & 捕捉异常时段行为与时序突变\\
    账户标识及账户历史 & 账户时间窗交易次数、平均金额、活跃度统计 & 类别索引 + 统计型 & 安全编码、滑动窗口聚合 & 表征账户级长期与短期行为模式\\
    设备标识 & 设备使用频次、设备共享度、账户---设备切换统计 & 类别型/统计型 & 缺失标识、频次编码、窗口统计 & 识别设备复用、盗刷与团伙风险\\
    商户标识/商户类别 & 商户类别编码、商户频次、商户局部风险统计 & 类别型/统计型 & one-hot、频次编码或安全统计编码 & 增强交易场景相似性与商户异常识别能力\\
    地理位置字段 & 地域编码、跨地域切换次数、地理跳变标记 & 类别型/统计型 & 地域粒度统一、编码、窗口切换统计 & 捕捉异地交易、短时跨区域异常\\
    支付渠道/终端类型 & 渠道编码、终端类别向量、渠道频次统计 & 类别型 & one-hot 或频次编码 & 丰富候选召回的相似度基础，提升异常区分度\\
    图节点原始属性 & 节点属性规范化向量 & 混合型 & 数值归一化、类别编码 & 形成图节点输入的基础表示\\
    图边关系与邻域信息 & 入边/出边计数、邻域金额统计、方向比例、时间分布 & 关系统计型 & 邻域聚合、时间窗统计、向量拼接 & 将图结构信息映射为可检索、可评分的节点向量\\
    % TODO(author): 运行最终预处理脚本输出具体 d，并同步回填第三章表 3-2。
    融合后统一表示 & 固定维度最终向量 ($v_i\in\mathbb{R}^d$，$d$ 由最终预处理脚本确定) & 向量型 & 维度检查、类型检查、异常值截断 & 作为 FlashANNS/FAISS-GPU 候选召回输入，并直接供 PyTOD 评分使用\\
    \bottomrule
  \end{tabular}
  \end{adjustbox}
\end{table}

\section{多 GPU 实验配置补充表}

\begin{table}[htbp]
  \caption{多 GPU 实验配置补充表}
  \label{tab:multi-gpu-exp-setup}
  \centering
  \small
  \setlength{\tabcolsep}{3pt}
  \renewcommand{\arraystretch}{1.18}
  \textbf{（a）固定参数与实验口径}
  \begin{adjustbox}{max width=\linewidth}
    \begin{tabular}{C{0.9cm}P{3.1cm}C{1.2cm}C{1.2cm}C{1.2cm}C{0.8cm}P{4.3cm}}
      \toprule
      GPU数 & 方案 & 数据规模 & 批大小 & 候选数 & $k$ & 实验目的 \\
      \midrule
      1 & Replicated & 10M & 4096 & 128 & 50 & 单卡吞吐扩展基线 \\
      2 & 索引复制 & 10M & 4096 & 128 & 50 & 2 GPU 索引复制扩展基线 \\
      4 & 索引复制 & 10M & 4096 & 128 & 50 & 四卡索引复制扩展基线 \\
      4 & 本地评分 + 设备侧归并 & 100M & 3072 & 128 & 50 & 显存受限时的设备侧归并兜底方案 \\
      4 & 本地评分 + 主机侧聚合 & 100M & 3072 & 128 & 50 & 主机侧聚合对比基线 \\
      2 & 拓扑感知绑定 & 100M & 4096 & 128 & 50 & I/O 通道绑定对比实验 \\
      4 & 拓扑感知绑定 & 100M & 4096 & 128 & 50 & I/O 通道绑定对比实验 \\
      2 & Shared-path conflict & 100M & 4096 & 128 & 50 & 共享路径冲突对比基线 \\
      4 & Shared-path conflict & 100M & 4096 & 128 & 50 & 共享路径冲突对比基线 \\
      \bottomrule
    \end{tabular}
  \end{adjustbox}

  \vspace{0.6em}
  \textbf{（b）方案配置说明}
  \begin{adjustbox}{max width=\linewidth}
    \begin{tabular}{P{3.2cm}P{3.0cm}P{3.5cm}P{3.1cm}P{3.4cm}}
      \toprule
      方案 & 索引组织 & 归并路径 & I/O 绑定 & 定位 \\
      \midrule
      Replicated & Full replica & None & Topology-aware & 主扩展方案，验证随 GPU 数量增加的吞吐收益 \\
      本地评分 + 设备侧归并 & 分片索引 & NCCL all-gather + 设备侧归并 & 拓扑感知 & 显存受限时的设备侧归并兜底路径 \\
      本地评分 + 主机侧聚合 & 分片索引 & 主机侧聚合 & 拓扑感知 & 用于对比 CPU 回到数据平面的额外代价 \\
      拓扑感知绑定 & 完整副本 & 无 & 拓扑感知 & 验证 I/O 通道绑定对扩展效率的影响 \\
      Shared-path conflict & Full replica & None & Shared PCIe/SSD path & 作为共享路径冲突下的对比基线 \\
      \midrule
      \multicolumn{5}{p{14.8cm}}{\footnotesize 注：多 GPU 实验规模口径统一按“真实集 $\rightarrow$ 1M 合成数据 $\rightarrow$ 10M 合成数据 $\rightarrow$ 50M 合成数据 $\rightarrow$ 100M 合成数据”组织。10M 主要服务于索引复制主方案扩展曲线，100M 主要服务于兜底归并与路径冲突对比。} \\
      \bottomrule
    \end{tabular}
  \end{adjustbox}
\end{table}

\chapter{缩写、变量与配置对照}

本附录列出正文中频繁出现的关键变量、配置项及其含义，便于后续统一实验脚本和结果复核。

\renewcommand{\arraystretch}{1.2}
\small
\begin{tabularx}{\textwidth}{P{0.34\textwidth}Y}
变量或配置项 & 含义 \\
\texttt{batch\_size} & 单次提交到候选召回或精排路径中的批处理大小。 \\
\texttt{rerank\_budget} & 进入高精度精排阶段的候选数量预算。 \\
\texttt{nlist} & 索引聚类桶数量或粗排分区规模参数。 \\
\texttt{nprobe} & 查询阶段实际访问的粗排桶数量。 \\
\texttt{device\_id} & 当前任务绑定的 GPU 设备标识。 \\
\texttt{replica\_mode} & 多 GPU 复制式多实例执行模式开关。 \\
\texttt{bounded\_batch\_size} & 在时延约束下允许形成的微批上限。 \\
\end{tabularx}
